\documentclass[11pt]{article}
\usepackage[a4paper,margin=1in]{geometry}
\usepackage[T1]{fontenc}
\usepackage{enumitem}
\usepackage{booktabs}
\usepackage{listings}
\usepackage{xcolor}
\usepackage{parskip}
\usepackage[hidelinks]{hyperref}

\definecolor{codebg}{gray}{0.96}
\lstset{
  basicstyle=\ttfamily\small,
  backgroundcolor=\color{codebg},
  frame=single,
  framerule=0pt,
  breaklines=true,
  columns=fullflexible,
  keepspaces=true,
  xleftmargin=6pt,
  xrightmargin=6pt,
  aboveskip=8pt,
  belowskip=8pt,
}

\title{\textbf{Building \texttt{BarsN} and \texttt{BarsP} from source\\
\large An operating-system--independent guide}}
\author{}
\date{\today}

\begin{document}
\maketitle

\section{The one idea to remember}

\texttt{BARS} ships C source that must be compiled into two products per method:
a stand-alone \emph{executable} (\texttt{barsN.out}, \texttt{barsP.out}) and a
\emph{shared library} (\texttt{barsN.so}, \texttt{barsP.so}). Compiled files are
tied to one operating system \emph{and} one CPU architecture. The portable rule,
true on every platform, is therefore:

\begin{quote}
\textbf{Copy the sources, not the binaries. Rebuild on each machine.}
\end{quote}

Never reuse a \texttt{*.o}, \texttt{*.so}, or \texttt{*.out} produced on a
different machine: an object of the wrong architecture loaded by R is exactly the
``incompatible architecture'' / ``found $\ldots$ required $\ldots$'' error. The
same source tree, compiled locally, works everywhere.

\section{What is platform-specific and what is not}

Almost every correction the code needs is about \emph{modern C compilers and
portability}, not about a particular CPU or OS. Only two things truly vary by
platform: the toolchain you install, and the flags that name the system BLAS/LAPACK.

\begin{center}
\begin{tabular}{@{}lll@{}}
\toprule
\textbf{Aspect} & \textbf{Same everywhere?} & \textbf{Note} \\
\midrule
\texttt{malloc.h}$\rightarrow$\texttt{stdlib.h} edit & yes & portability fix \\
Relaxed compiler flags        & yes & needed by clang 15+ / gcc 14+ \\
Two build products (.out, .so)& yes & \texttt{make} + \texttt{R CMD SHLIB} \\
File layout (exec/, libs/)    & yes & package convention \\
\midrule
C compiler                    & \textbf{no} & clang / gcc / Rtools \\
BLAS \& LAPACK link flags      & \textbf{no} & resolved via \texttt{R CMD config} \\
CPU architecture of output    & \textbf{no} & must match your R build \\
\bottomrule
\end{tabular}
\end{center}

The trick that makes the build portable is to stop hard-coding library names and
let \textbf{R itself} report the correct BLAS/LAPACK/compiler for the current
machine (Section~\ref{sec:build}).

\section{Step 0 --- Prerequisites}

You need a C compiler and a build of R whose architecture matches the machine.

\begin{itemize}[leftmargin=1.4em]
  \item \textbf{macOS:} Xcode Command Line Tools --- \texttt{xcode-select --install}
        (provides \texttt{clang}).
  \item \textbf{Linux:} the compiler and R headers, e.g.\ Debian/Ubuntu
        \texttt{sudo apt-get install build-essential r-base-dev}; Fedora/RHEL
        \texttt{sudo dnf install gcc make R-devel}.
  \item \textbf{Windows:} \textbf{Rtools} matching your R version, on the build PATH.
\end{itemize}

Confirm the compiler and R agree on architecture before building:

\begin{lstlisting}
cc --version        # or: gcc --version
R --version         # architecture here must match the compiler
\end{lstlisting}

On Apple Silicon, \texttt{R --version} must report \texttt{aarch64}; if it says
\texttt{x86\_64} you are running the Intel build under Rosetta and the output will
be built for the wrong architecture.

\section{Step 1 --- Bring the corrected sources across}

Copy the \texttt{barsN} and \texttt{barsP} source folders. \textbf{First remove any
compiled products that travelled with them}, since they belong to another machine:

\begin{lstlisting}
cd barsN && rm -f *.o *.so *.out
cd ../barsP && rm -f *.o *.so *.out
\end{lstlisting}

If you are starting from pristine upstream sources, apply the one portability edit
that BARS needs --- \texttt{<malloc.h>} does not exist on all systems, and
\texttt{<stdlib.h>} provides the same declarations portably:

\begin{lstlisting}
# in both barsN/ and barsP/, in logspline.c (line 3):
#   change   #include <malloc.h>
#   to       #include <stdlib.h>
\end{lstlisting}

This edit is correct on every operating system (it is required on macOS and
harmless on Linux and Windows).

\section{Step 2 --- Build both products, portably}
\label{sec:build}

Two separate builds are required for each method:
\begin{itemize}[leftmargin=1.4em,itemsep=2pt]
  \item the \textbf{executable} \texttt{barsN.out} --- this is what the package's
        R code actually runs, via \texttt{system()};
  \item the \textbf{shared library} \texttt{barsN.so} --- placed in the package's
        \texttt{libs/} directory per its layout.
\end{itemize}

Rather than hard-code \texttt{-llapack -lblas} (and the obsolete \texttt{-lg2c},
which no modern system has), we ask R for the flags it was built with. These
commands are \textbf{identical on macOS, Linux, and Windows}; only R's answers
differ underneath.

\subsection*{BarsN}

\begin{lstlisting}
cd barsN
CC=$(R CMD config CC)
LIBS="$(R CMD config LAPACK_LIBS) $(R CMD config BLAS_LIBS) $(R CMD config FLIBS)"
FLAGS="-std=gnu89 -Wno-implicit-function-declaration -Wno-implicit-int"

# (a) the executable
$CC $FLAGS -o barsN.out \
    barsN.c barsN_funcs.c barsN_utils.c logspline.c ranlibsub.c com.c $LIBS

# (b) the shared library
PKG_CFLAGS="$FLAGS" PKG_LIBS="$LIBS" \
R CMD SHLIB -o barsN.so \
    barsN.c barsN_funcs.c barsN_utils.c logspline.c ranlibsub.c com.c
\end{lstlisting}

\subsection*{BarsP}

Identical, with \texttt{P} for \texttt{N} (ensure the \texttt{malloc.h} edit is
present first):

\begin{lstlisting}
cd ../barsP
$CC $FLAGS -o barsP.out \
    barsP.c barsP_funcs.c barsP_utils.c logspline.c ranlibsub.c com.c $LIBS

PKG_CFLAGS="$FLAGS" PKG_LIBS="$LIBS" \
R CMD SHLIB -o barsP.so \
    barsP.c barsP_funcs.c barsP_utils.c logspline.c ranlibsub.c com.c
\end{lstlisting}

The \texttt{-std=gnu89} and \texttt{-Wno-implicit-*} flags let this legacy
C89 code compile under modern clang and gcc, which now treat implicit
declarations and implicit \texttt{int} as errors by default. Expect only harmless
\texttt{-Wformat-security} and unused-variable warnings.

\paragraph{What \texttt{R CMD config} resolves to.}
On macOS this expands to Apple's \texttt{Accelerate} framework; on Linux to R's
own \texttt{libRblas}/\texttt{liblapack}; on Windows to the libraries bundled with
Rtools. You do not need to know which --- that is the point of asking R.

\section{Step 3 --- Confirm the architecture}

The check that matters most: executable, library, and R must all be the same
architecture.

\begin{lstlisting}
file barsN.out barsN.so barsP.out barsP.so
\end{lstlisting}

On macOS each line should report \texttt{arm64} (Apple Silicon) or
\texttt{x86\_64} (Intel), matching your R; on Linux \texttt{x86-64} or
\texttt{aarch64}. If any product reports the wrong architecture it is a foreign
copy that slipped through: \texttt{rm -f *.o *.so *.out} and rebuild.

\section{Step 4 --- Place the products in the package}

Put the executables in \texttt{exec/} and the shared libraries in \texttt{libs/}:

\begin{lstlisting}
cp barsN/barsN.out barsP/barsP.out   BARS/exec/
cp barsN/barsN.so  barsP/barsP.so    BARS/libs/
\end{lstlisting}

The package is then ready to install (or re-package as a source tarball with
\texttt{R CMD build BARS}; remember such a tarball should carry \emph{source},
so each target machine rebuilds its own binaries).

\section{The \texttt{\textasciitilde/.R/Makevars} caveat}

If you carry over a personal \texttt{\textasciitilde/.R/Makevars} from another
machine, any line naming a machine-specific compiler or library path (a
hard-coded \texttt{gfortran} location in \texttt{FLIBS}, for instance) will be
wrong on the new machine and can derail unrelated source installs. Either delete
such lines and let R use its own defaults, or repoint them at the new machine's
toolchain. Architecture-independent definitions (such as
\texttt{-DCalloc=R\_Calloc -DFree=R\_Free -DRealloc=R\_Realloc}) are safe to keep.

\section*{One-line summary}

\begin{quote}
Copy only the C sources; apply the \texttt{malloc.h}$\rightarrow$\texttt{stdlib.h}
edit; build \texttt{.out} with the compiler and \texttt{.so} with
\texttt{R CMD SHLIB}, taking BLAS/LAPACK flags from \texttt{R CMD config}; verify
with \texttt{file}; then drop \texttt{.out} into \texttt{exec/} and \texttt{.so}
into \texttt{libs/}. The procedure is the same on every operating system --- only
the toolchain you installed differs.
\end{quote}

\end{document}
